\documentclass[10pt,twocolumn]{article}
\usepackage[letterpaper,margin=0.68in,columnsep=0.24in]{geometry}
\usepackage{amsmath,amssymb,amsthm,booktabs,microtype,xcolor}
\usepackage[hidelinks]{hyperref}
\usepackage{enumitem}
\setlist{nosep,leftmargin=*}
\newtheorem{proposition}{Proposition}
\newtheorem{corollary}{Corollary}
\newtheorem{remark}{Remark}
\newcommand{\R}{\mathbb{R}}
\newcommand{\E}{\mathbb{E}}
\title{\textbf{The Ecopoietic Hedge: When Symbiont Diversity Improves Restoration Reliability Under Environmental Uncertainty}}
\author{Justin D. Hart\\Viridis Research\\\texttt{viridisnorthllc@gmail.com}}
\date{2 August 2026}

\begin{document}
\maketitle

\begin{abstract}
Microbial and mycorrhizal inocula can support plant establishment, but their
effects depend on host, stress, delivery practice, and environmental context.
This creates a decision problem upstream of any field deployment: should a
restoration program allocate a fixed inoculum dose to the strain with the best
point estimate, or hedge across strains before the realized establishment
regime is known? We analyze a preregistered two-strain, two-regime model in
which effective exposures are additive in the negative log failure
probability. When strain performance crosses symmetrically across regimes,
with efficacy matrix $[(\alpha,\beta),(\beta,\alpha)]$ and
$\alpha>\beta>0$, the equal-dose mixture uniquely maximizes worst-regime
establishment. Its guaranteed log-failure exponent is
$(\alpha+\beta)D/2$, compared with $\beta D$ for either monoculture. For a
fixed target probability, the implied dose saving is
$(\alpha-\beta)/(\alpha+\beta)$. The conclusion has sharp negative
controls: identical strains provide no diversification benefit, and a strain
that weakly dominates in every modeled regime should receive the entire dose.
An exact breakpoint solver, grid-convergence tests, 1,000 randomized symmetric
instances, 250 multi-scenario instances, and numerical-stability tests passed
32 of 32 preregistered checks. A separate Lean 4 artifact proves the five
load-bearing model claims and a concrete non-vacuity witness with no remaining
proof holes. These are synthetic model and formal checks, not field
verification. The result supplies a falsifiable criterion for when a
symbiont portfolio is justified: measured cross-regime rank reversal, not
diversity by itself.
\end{abstract}

\noindent\textbf{Keywords:} ecological restoration; plant microbiome;
symbiosis; robust optimization; microbial inoculants; environmental
uncertainty; maximin allocation

\section{Introduction}

Large-scale restoration increasingly treats microbial partners as active
components of establishment rather than invisible background. Mycorrhizal
fungi can alter nutrient and water acquisition, plant-growth-promoting
bacteria can change stress physiology, and designed synthetic communities can
carry several complementary functions. Yet biological efficacy is not a
property of a strain alone. It depends on the host, soil, drought intensity,
timing, community history, and delivery method. Current experiments make this
context dependence concrete. Wang et al. found that selected
\emph{Rhizobium} strains and synthetic communities improved Arabidopsis growth
under osmotic stress while having little effect under the mock condition
\cite{wang2024}. Bandopadhyay et al. separated plant- and environment-mediated
drivers of active rhizosphere responses to drought and found no single taxon
that behaved as a drought generalist across their two crop systems
\cite{bandopadhyay2024}. Lin et al. showed that inoculation practice affected
early fabricated-community assembly and that reproducibility differed across
delivery protocols \cite{lin2024}.

Combined inocula can help, but not uniformly. In a Mediterranean maize field
experiment, single and dual microbial treatments changed distinct plant and
soil outcomes under watered and drought conditions \cite{ouhaddou2024}. A
16-member drought-tolerant bacterial synthetic community improved barley
performance in a pot experiment \cite{rigerte2025}. Conversely, field
mycorrhizal inoculation during drought aided one forage species while
increasing competition and water stress for another \cite{amf2023}. These
results motivate portfolios without proving that more strains are always
better. A mixture consumes finite dose, can include redundant organisms, and
can introduce antagonism or biosecurity risk.

The decision studied here is deliberately narrower than biological community
design in full. A manager has two approved, non-antagonistic candidate strains,
a fixed total dose, and multiple plausible establishment regimes. Allocation
must occur before the realized regime is known. We ask when the mixture raises
the guaranteed probability of at least one effective exposure.

This question is distinct from two nearby results in the Viridis working
corpus. Run 095 studied the order of one-way facilitation across biosphere
layers \cite{run095}; Run 096 studied deterministic plant-symbiont
co-establishment and best-strain matching for a known host-site pair
\cite{run096}. Here the host and total dose are fixed. The new axis is
pre-regime robust allocation across alternative symbionts. The work remains a
generated exploration and novelty is unassessed. Its five preregistered formal
targets were subsequently proved in Lean 4 and independently rebuilt and
audited after Aristotle review. Formal verification establishes only the
stated mathematics; it does not establish the ecological assumptions or field
efficacy.

\section{Prior work and decision context}

The biological literature contains both reasons to diversify and reasons to
avoid an automatic diversity rule. Synthetic communities can combine traits
that no single strain carries, but community assembly is sensitive to
inoculation practice \cite{lin2024}. Host genotype and root traits can recruit
different microbes under stress \cite{wang2024}; environmental exposure can
independently restructure the active community \cite{bandopadhyay2024}. A
strain portfolio is therefore best interpreted as a hedge over unresolved
host--environment response, not as a universal ecological good.

The mathematical form is a small robust optimization problem. Robust
optimization selects a decision that remains feasible or performant across a
specified uncertainty set rather than optimizing only one point estimate
\cite{bental1998}. In this paper the uncertainty set is finite and explicit:
each scenario supplies one efficacy per strain. The objective is maximin, so it
prioritizes the weakest modeled regime. This is appropriate only when the
scenarios are safety-relevant and no defensible probability distribution is
available. If credible probabilities exist, expected utility or a coherent
risk measure may be preferable.

We model establishment through a Poisson-exposure approximation. If an
effective colonization event occurs with exponent $E\ge0$, then the probability
of at least one event is
\begin{equation}
 p(E)=1-\exp(-E).
 \label{eq:success}
\end{equation}
The exponent is additive across independent doses. Equation
\eqref{eq:success} is an operational approximation, not a biological law. Its
advantage is that dose contributions combine in negative log failure, allowing
the allocation question to be stated exactly and tested directly.

\section{Model, hypotheses, and failure conditions}

Let $D>0$ be total dose. Allocate $x\in[0,D]$ to strain A and $D-x$ to strain B.
There are two plausible regimes. In regime 1, efficacies per dose are
$(\alpha,\beta)$; in regime 2 they are $(\beta,\alpha)$, with
$\alpha>\beta>0$. The regime exponents are
\begin{align}
 E_1(x)&=\alpha x+\beta(D-x)
       =\beta D+(\alpha-\beta)x,\\
 E_2(x)&=\beta x+\alpha(D-x)
       =\alpha D-(\alpha-\beta)x.
\end{align}
Because $p(E)$ is strictly increasing, maximizing worst-regime success is
equivalent to
\begin{equation}
 \max_{0\le x\le D} g(x),\qquad
 g(x)=\min\{E_1(x),E_2(x)\}.
 \label{eq:maximin}
\end{equation}

The hypothesis, assumptions, and stopping rules were written into
\texttt{verification.py} before this manuscript. The assumptions are:
\begin{enumerate}
\item log-failure contributions add within each regime;
\item per-dose costs are equal and total dose is fixed;
\item allocation precedes regime realization and uses a worst-case objective;
\item efficacies are nonnegative and fixed over the establishment window; and
\item the symmetric matrix is a test case, not a measured site claim.
\end{enumerate}

The preregistered claim would be abandoned if dense search moved the optimum
away from $D/2$ beyond grid resolution, if the analytic value disagreed with an
independent breakpoint solver, if identical strains retained a claimed gain,
if mixing beat a uniformly dominant strain, or if refinement failed to
converge.

\section{Analytic results}

\begin{proposition}[Symmetric robust allocation]
\label{prop:equal}
For $\alpha>\beta>0$ and $D>0$, the unique maximizer of
\eqref{eq:maximin} is $x^*=D/2$, and
\begin{equation}
 g(x^*)=\frac{\alpha+\beta}{2}D.
 \label{eq:star}
\end{equation}
\end{proposition}

\begin{proof}
$E_1$ is strictly increasing in $x$ and $E_2$ is strictly decreasing. Their
unique crossing solves
$\beta D+(\alpha-\beta)x=\alpha D-(\alpha-\beta)x$, hence $x=D/2$.
To the left of the crossing, $E_1<E_2$, so $g=E_1$ increases strictly. To the
right, $E_2<E_1$, so $g=E_2$ decreases strictly. The crossing is therefore the
unique maximizer. Substitution gives \eqref{eq:star}.
\end{proof}

\begin{corollary}[Guaranteed exponent gain]
\label{cor:gain}
Either monoculture has worst-regime exponent $\beta D$. The mixed allocation
improves that exponent by
\begin{equation}
 \Delta g=\frac{\alpha-\beta}{2}D>0.
\end{equation}
\end{corollary}

\begin{corollary}[Target-probability dose]
\label{cor:dose}
For target success $p_0\in(0,1)$, define $L=-\log(1-p_0)$. The robust doses are
\begin{equation}
 D_{\rm mix}=\frac{2L}{\alpha+\beta},\qquad
 D_{\rm mono}=\frac{L}{\beta}.
\end{equation}
The fractional dose saving is
\begin{equation}
 1-\frac{D_{\rm mix}}{D_{\rm mono}}
 =\frac{\alpha-\beta}{\alpha+\beta}.
 \label{eq:saving}
\end{equation}
\end{corollary}

The result is not a generic endorsement of equal splitting. Consider $m$
scenarios with efficacy pairs $(a_s,b_s)$. The two-strain robust problem is
\begin{align}
 \max_{x,t}\quad &t\\
 \text{subject to}\quad&a_s x+b_s(D-x)\ge t,\quad s=1,\ldots,m,\\
 &0\le x\le D.
 \label{eq:lp}
\end{align}
The lower envelope of the scenario lines is concave and piecewise linear, so an
optimum occurs at an endpoint or at an intersection of two scenario lines.
This supplies an exact breakpoint algorithm for two strains. In an asymmetric
crossover, the robust solution can be interior without being $D/2$.

\begin{proposition}[Dominance negative control]
\label{prop:dominance}
If $a_s\ge b_s$ for every modeled scenario, $x=D$ is maximin. If the inequality
is strict in every active worst-case scenario, moving dose from A to B strictly
reduces the robust objective locally.
\end{proposition}

\begin{proof}
For any $x\le D$ and each $s$,
$a_sD-[a_sx+b_s(D-x)]=(a_s-b_s)(D-x)\ge0$. Thus the all-A allocation weakly
improves every scenario simultaneously and cannot have a smaller minimum.
\end{proof}

If $a_s=b_s$ for every scenario, each line is constant in $x$. Every allocation
ties, so merely increasing strain count has no modeled value. These controls
are decision-relevant: the portfolio is justified by crossing performance, not
by a diversity label.

\section{Numerical verification}

The verification program was frozen before the paper and uses only the Python
standard library. It independently implements the formulas, an exact
two-strain breakpoint solver, and a dense-grid solver. All 32 checks passed.

For the illustrative values $\alpha=1.7$, $\beta=0.3$, and $D=4$, the equal
mix has exponent $4$ and worst-regime success
$1-e^{-4}=0.98168436$. Either monoculture has exponent $1.2$ and success
$1-e^{-1.2}=0.69880579$. The implied target-dose saving is $0.7$. These are
dimensionless synthetic values, not treatment recommendations.

Three odd grid resolutions---101, 1,001, and 10,001 intervals---were chosen so
that the equal split was not sampled exactly. Value errors fell from
$2.772\times10^{-2}$ to $2.797\times10^{-3}$ to
$2.800\times10^{-4}$, within the derived slope-resolution bound. The exact
solver matched the closed form over 1,000 randomized symmetric parameter sets;
the largest allocation discrepancy was $6.253\times10^{-13}$ and the largest
value discrepancy was $5.684\times10^{-14}$.

For 250 random problems with two to seven scenario lines, the exact breakpoint
solution agreed with a 20,000-interval grid within the relevant slope bound;
the largest observed value gap was $1.704\times10^{-4}$. The identical-strain
control was flat. In four dominance scenarios, the exact solver assigned all
dose to the dominant strain. An asymmetric three-scenario crossover produced
an interior allocation $x=1.5$ for $D=4$, explicitly disproving the use of
50/50 as a general heuristic.

The probability path uses \texttt{expm1} to avoid cancellation at small
exponents. A value of $10^{-14}$ remained resolvable, while an exponent of
1,000 saturated safely at probability one. The complete output is retained in
\texttt{verification\_output.txt}; no test output was truncated.

The correct numerical status is \texttt{NUMERICAL\_GATE\_PASS}. The program
checks a finite implementation and does not by itself provide Lean proof,
field validation, novelty assessment, or evidence of ecological safety. The
separate formal artifact described next addresses the bounded mathematical
claims only.

\section{Formal verification}

The audited Lean 4 module \path{EcopoieticHedge.lean} uses toolchain
\path{leanprover/lean4:v4.28.0} and the pinned mathlib revision recorded in the
accompanying forge request. It contains five load-bearing theorems plus one
explicit non-vacuity witness:
\begin{enumerate}
\item \path{reciprocal_two_regime_maximin_equal_split};
\item \path{reciprocal_diversification_gain};
\item \path{target_probability_dose_saving};
\item \path{uniform_dominance_implies_monoculture_optimal};
\item \path{identical_strains_zero_diversification_gain}; and
\item \path{ecopoietic_hedge_nonvacuous}.
\end{enumerate}
The first theorem proves both the maximin upper bound and the equality
condition $x=D/2$. The second states the exact exponent gain against the
$x=0$ monoculture endpoint; direct evaluation of the reciprocal definitions
gives the same endpoint value at $x=D$. The third retains
$\alpha>\beta>0$ and $0<p<1$. The fourth proves that the all-A allocation
weakly improves every scenario under pointwise dominance, which is the
sufficient maximin certificate used in Proposition~\ref{prop:dominance}. The
fifth proves allocation-independence under identical scenario coefficients.
The witness evaluates $\alpha=2$, $\beta=1$, and $D=2$, obtaining mixed and
monoculture worst-case exponents 3 and 2.

The local forge audit rebuilt the returned project, compared every submitted
definition and theorem signature with the frozen request, scanned the Lean
sources for proof holes and prohibited escapes, inspected the axioms of all
six named theorems, checked the explicit witness, and reconciled the target
names against the paper claim inventory. This verifies the formal model only;
the additive exposure law, parameter relevance, and biological interpretation
remain assumptions to be tested empirically.

\section{Predictions and empirical tests}

The central empirical prediction is conditional and discriminating:

\begin{quote}
An inoculant portfolio should improve held-out worst-regime establishment over
the best monoculture only when estimated strain performance crosses across the
relevant regimes; under uniform dominance, the portfolio should not help after
holding total viable dose fixed.
\end{quote}

A direct test is a host-by-strain-by-regime factorial experiment with at least
four dose levels plus an uninoculated control. Regimes should be chosen before
outcome inspection and represent operational uncertainty such as watered versus
droughted soil, low versus high phosphorus, or two delivery timings. Within
each cell, estimate establishment failure $q$ and test whether
$-\log q$ is approximately linear in viable dose over the proposed operating
range. The strain efficacy is the slope, with uncertainty propagated from
replicate blocks.

The efficacy matrix should then be classified before portfolio evaluation:
\begin{enumerate}
\item \textbf{crossover:} rank order reverses across at least two regimes;
\item \textbf{dominance:} one strain is no worse in every regime; or
\item \textbf{unresolved:} confidence intervals do not distinguish either.
\end{enumerate}
For a crossover matrix, freeze the robust allocation and compare it with both
monocultures on held-out blocks. Refutation occurs if the additive hazard model
fails materially, if antagonistic interactions make the combined exponent
non-additive, or if the predicted mixture fails to improve held-out worst-
regime establishment. A dominance result favoring monoculture confirms the
negative control rather than invalidating the framework.

Secondary tests should estimate interaction terms by comparing the mixture
with the sum of single-strain log-failure exponents. A negative interaction
large enough to reverse the robust gain is a stop condition. Long-term
persistence, ecosystem composition, pathogen transmission, and off-target
effects must be measured separately; early establishment is not restoration
success.

\section{Discussion}

The mathematical result turns a vague case for microbial diversity into an
auditable decision rule. A mix is useful when it spans environmental response
directions that a single strain does not. In the symmetric case, the two
regime constraints pull allocation in opposite directions and the robust
solution equalizes them. The same geometry appears in reliability portfolios:
diversification is valuable when components fail differently, not when they
are clones.

This clarifies the relationship to deterministic symbiont matching. If the
site regime is known and one strain has the highest efficacy, Run 096's
matching logic remains appropriate. The hedge becomes relevant only when the
decision precedes unresolved regime realization and the cost of failure makes
the lower tail load-bearing. It therefore extends, rather than replaces, the
deterministic pair model.

The result also places a burden on data collection. Without cross-regime dose
response, a portfolio recommendation is not evidence based. A single-condition
screen can rank strains but cannot estimate hedge value. Conversely, a small
factorial trial can decide that the mix is unnecessary, which is economically
useful because it avoids spending dose on a dominated strain.

No claim is made that the current primary literature contains the same
closed-form model or that the result is new to robust optimization, ecological
insurance theory, or inoculant design. The literature search was targeted, not
systematic. Novelty remains unassessed pending independent prior-art and
significance review.

\section{Limitations}

The model omits most of the ecology that determines deployment safety and
durability. Independent Poisson exposures can fail when strains compete,
facilitate one another, share a delivery bottleneck, or change host recruitment.
Efficacy can vary with dose, time, genotype, soil legacy, and resident community.
Equal per-dose cost ignores production yield, shelf life, regulatory burden,
and viable propagule uncertainty. The maximin objective can be overly
conservative, and the two-regime symmetric case is intentionally stylized.

Early establishment probability is not ecosystem function, persistence,
native-community recovery, or net conservation benefit. A portfolio can also
increase biosecurity and pathogen risks. Only approved, locally appropriate
organisms should enter a field trial, under the applicable ecological and
regulatory review. The numerical gate does not evaluate any of those conditions.

The general two-strain finite-scenario problem is easy; many-strain robust
allocation with interactions, integer packaging, capacity limits, uncertain
efficacies, and chance constraints is materially harder. Those extensions are
deferred. So are empirical calibration and comparison with simpler baselines.

\section{Reproducibility}

\paragraph{Reproducibility.}
Run \texttt{python3 verification.py} with Python 3. The program uses only the
standard library and a fixed pseudorandom seed, prints every aggregate check,
and exits nonzero on failure. The exact source and complete output are included
with this paper. The manuscript can be rebuilt with two passes of
\texttt{pdflatex -interaction=nonstopmode -halt-on-error paper.tex}.

\paragraph{Data and code availability.}
No external dataset was analyzed. All generated code, synthetic parameters,
claims, and captured output needed to reproduce the numerical result are in the
immutable Run 117 package. External papers are identified in
\texttt{SOURCE\_LEDGER.json} by DOI or stable URL.

The formal artifact is rebuilt with \texttt{lake build} using the included
\texttt{lean-toolchain}, \texttt{lakefile.toml}, and pinned
\texttt{lake-manifest.json}. The immutable Aristotle requests, returned Lean
tree, repair receipt, and seven-gate forge audit are included in the staged
submission package.

\paragraph{Accountable authorship.}
Justin D. Hart is the accountable human author and is responsible for the
scientific claims, source selection, release decisions, and any future
correction. No AI system is an author.

\paragraph{Substantive AI-use disclosure.}
OpenAI Codex generated the initial hypothesis, performed the local corpus and
literature search under Justin D. Hart's standing direction, wrote the
preregistered verification program, executed the synthetic checks, drafted the
manuscript and package metadata, and checked compilation and rendering. Codex
also narrowed the claim to a conditional crossover result and retained the
negative controls. Historical Claude-generated Viridis artifacts were read as
attributed prior work and were not rewritten. Aristotle was not used for this
run's initial hypothesis, source work, numerical program, or manuscript. It
was used as the external formal-verification service for the hash-bound Lean
request and repair review. Codex developed and locally compiled the bounded
proof repair after the first Aristotle task stalled before proof work, then
performed the separate deterministic statement, build, escape, axiom,
non-vacuity, and significance audit. The human author remains accountable for
review and any external use.

\paragraph{Funding.}
No external funding was received specifically for this generated exploration.

\paragraph{Conflicts of interest.}
The author is the founder of Viridis and could benefit if a future decision
tool based on this work is commercialized. No signed revenue, customer
validation, adoption, or product efficacy is evidenced by this paper.

\section{Conclusion}

Under a transparent additive log-failure model, symbiont diversity is a robust
hedge only when strain performance crosses across unresolved establishment
regimes. Symmetric crossover yields an equal split and an exact dose-saving
formula; identical strains yield no gain; uniform dominance favors a
monoculture. The strongest next step is not a broader claim but a small
preregistered factorial experiment that estimates the efficacy matrix and
tries to falsify additivity on held-out blocks. The bounded model claims now
have zero-hole Lean proofs, but the result remains a synthetic generated
exploration, not field guidance, empirical validation, or a novelty claim.

\begin{thebibliography}{99}

\bibitem{wang2024}
Z. Wang, Z. Li, Y. Zhang, et al., ``Root hair developmental regulators
orchestrate drought triggered microbiome changes and the interaction with
beneficial Rhizobiaceae,'' \emph{Nature Communications} \textbf{15}, 10068
(2024). doi:10.1038/s41467-024-54417-5.

\bibitem{bandopadhyay2024}
S. Bandopadhyay, X. Li, A. W. Bowsher, R. L. Last, and A. Shade,
``Disentangling plant- and environment-mediated drivers of active rhizosphere
bacterial community dynamics during short-term drought,'' \emph{Nature
Communications} \textbf{15}, 6347 (2024).
doi:10.1038/s41467-024-50463-1.

\bibitem{lin2024}
H.-H. Lin, M. Torres, C. A. Adams, et al., ``Impact of Inoculation Practices on
Microbiota Assembly and Community Stability in a Fabricated Ecosystem,''
\emph{Phytobiomes Journal} \textbf{8}, 155--167 (2024).
doi:10.1094/PBIOMES-06-23-0050-R.

\bibitem{ouhaddou2024}
R. Ouhaddou, L. Ech-chatir, C. Ikan, et al., ``Investigation of the impact of
dual inoculations of arbuscular mycorrhizal fungi and plant growth-promoting
rhizobacteria on drought tolerance of maize grown in a compost-amended field
under Mediterranean conditions,'' \emph{Frontiers in Microbiology}
\textbf{15}, 1432637 (2024). doi:10.3389/fmicb.2024.1432637.

\bibitem{rigerte2025}
L. Rigerte, A. Heintz-Buschart, T. Reitz, and M. T. Tarkka, ``Assembly and
application of a synthetic bacterial community for enhancing barley tolerance
to drought,'' \emph{Frontiers in Bacteriology} \textbf{4}, 1572294 (2025).
doi:10.3389/fbrio.2025.1572294.

\bibitem{amf2023}
``Arbuscular mycorrhizal fungi increase the interspecific competition between
two forage plant species and stabilize the soil microbial network during a
drought event: Evidence from the field,'' \emph{Applied Soil Ecology}
\textbf{185}, 104805 (2023). doi:10.1016/j.apsoil.2023.104805.

\bibitem{bental1998}
A. Ben-Tal and A. Nemirovski, ``Robust convex optimization,''
\emph{Mathematics of Operations Research} \textbf{23}, 769--805 (1998).

\bibitem{run095}
J. D. Hart, ``The Ecopoietic Succession-Ordering Theorem,'' Viridis Nightly
Science Engine, Run 095 (2026). Internal working-corpus artifact.

\bibitem{run096}
J. D. Hart, ``The Mutualist Co-Establishment Theorem,'' Viridis Nightly Science
Engine, Run 096 (2026). Internal working-corpus artifact.

\end{thebibliography}

\end{document}
